\documentclass[11pt]{article}
\usepackage{icmlko}
\begin{document}

\title{한 자 차이로 서른 배 --- 두 해 묵은 \texttt{[:20]} 이 제일 값비싼 자료의 목을 조르고 있었다}
\author{월드모델 실험실}
\date{노트 491 · 2026년 8월 3일}
\maketitle

\begin{abstract}
노트 490 이 검색 시계열에 판 $-0.0283$ 을 매겼다 --- 두 해 동안 채택한
어떤 모형 변경보다 여섯 배 크다. 그런데 그 자료의 덮음은 아홉 도메인
합계 \textbf{$41.7\%$} 다. 열린 갈래는 두 해째 ``\emph{모바일 $\cdot$
게임 검색 덮음 $30 \sim 41\%$ 는 제목 언어 한계}''라고 적혀 있었다.

\textbf{언어가 아니었다.} 만화는 제목의 한글 비율이 \textbf{$0.000$}
(전부 라틴)인데 $5\%$ 가 붙고, 세계애니도 $0.000$ 인데 $18\%$ 가 붙는다.
게임은 \emph{붙는} 쪽이 오히려 라틴 비율이 $0.378$ 더 높다.

열 도메인 전부에서 일관된 신호는 하나뿐이었다 --- \textbf{길이}. 그리고
길이를 파고 들어가니 \texttt{ingest/trend\_all.py} 의 한 줄이 나왔다.

\begin{center}
\texttt{return (han[0] if han else segs[0])[:20]}
\end{center}

\textbf{조각을 고른 \emph{뒤에} 스무 자에서 자른다} --- 낱말 한가운데서.
``\emph{Family Go! - 대를 잇는 인생게임}''이 ``\emph{Family Go! - 대를
잇는 인}''이 되고, 데이터랩은 그런 질의에 아무것도 안 돌려준다.

\begin{center}
\begin{tabular}{lrr}
\hline
자르기 전 조각 길이 & 행 & 시계열이 붙음 \\
\hline
정확히 $20$ 자 (안 잘림) & $164$ & $\mathbf{18.3\%}$ \\
정확히 $21$ 자 (한 자 잘림) & $165$ & $\mathbf{0.6\%}$ \\
\hline
\end{tabular}
\end{center}

\textbf{한 자 차이에 서른 배}(피셔 $p = 1.1 \times 10^{-22}$ · 승산비
$29.2$).
\end{abstract}

\begin{figure}[t]
\centering
\includegraphics[width=\textwidth]{figs/f1.pdf}
\caption{\textbf{가} --- 자르기 전 길이별. \textbf{나} --- 경계 양옆.
\textbf{다} --- 되찾을 양.}
\end{figure}

\section{첫 읽기는 거의 틀렸다}

처음엔 \emph{질의 문자열}의 길이로 쟀다. 열 도메인 중 아홉에서 최대
길이가 정확히 $20$ 이고(만화 $60\%$ · 펀딩 $49\%$ · 세계애니 $45\%$ 가 그
$20$ 에 몰려 있다), $17 \sim 19$ 자는 $25.0\%$ 가 붙는데 $20$ 자는
$2.4\%$ 였다.

\textbf{그런데 그것만으로는 못 말한다.} 잘린 조각의 \emph{진짜} 길이는
중앙 $22 \sim 29$ 자다. 붙는 비율은 길이에 따라 원래 매끄럽게 내려가므로
($6$ 자 $62.9\%$ $\to$ $19$ 자 $23.6\%$), 진짜 길이 $25$ 자쯤이면
$2.4\%$ 는 \textbf{그 기울기의 연장일 수 있다}.

그래서 \textbf{자르기 전 길이}로 다시 그렸다. $20$ 자 미만은 매끄러운
로짓 직선이고(기울기 $-0.132$/자), 넘는 순간 그 직선 아래로 떨어진다 ---
$20$ 자 이상 전체가 실제 $2.8\%$ 대 예측 $11.2\%$ 로 \textbf{$-8.4\%$p}.
경계 바로 양옆($20$ 자 대 $21$ 자)에서 서른 배 차가 난다.

\begin{quote}
\textbf{조항} --- 절단을 의심할 때는 \textbf{자르기 전 값}으로 잰다.
자른 뒤 값으로 재면 절단의 효과와 원래 있던 기울기가 같은 축에 겹쳐
서로를 흉내 낸다.
\end{quote}

\section{되찾을 양}

$20$ 자 미만으로 맞춘 직선을 각 행의 진짜 길이에 넣어 예측하고 실제와의
차를 더하면 \textbf{$+152$ 행}이다 --- 만화 $+60$ · 세계애니 $+42$ ·
펀딩 $+20$ · 모바일 $+18$ · 애니 $+10$. 조인된 $6{,}420$ 행에서 덮음
$33.3\% \to 35.7\%$.

\textbf{예보지 결과가 아니다.} 다시 긁어야 확인된다. 고침은 넣었다 ---
한도는 남기되 \textbf{낱말 경계에서} 자른다(스무 자가 API 제약인지 우리
선택인지 기록이 없어, 어느 쪽이든 안전한 쪽으로).

\section*{남은 것}

\begin{itemize}
\item \textbf{다시 긁어야 한다} --- 이 노트는 진단이고, $+152$ 행은
예보다. 수집은 바깥으로 나가는 일이라 사람이 정한다.
\item \textbf{만화 $\cdot$ 세계애니의 진짜 벽은 따로 있다} --- 절단을
고쳐도 $9.1\%$ $\cdot$ $21.6\%$ 다. 한글이 한 글자도 없는 제목에
한국어 검색을 물어서다. \textbf{그건 언어 문제가 맞다} --- 다만
모바일 $\cdot$ 게임이 아니라 이 둘의 문제였다.
\item \textbf{두 해 동안 열린 갈래의 진단이 틀려 있었다.} 문장 하나가
``언어''라고 적혀 있어서 아무도 길이를 안 봤다. 갈래 목록의 \emph{설명}은
측정이 아니라 짐작이었는데 측정처럼 읽혔다.
\end{itemize}

\end{document}
